\section{Reinterpreting Renormalization}

\label{sec:renorm_reinterp}

The previous section redefined divergence
as contact with the \DH{}.

From this perspective,
the meaning of renormalization also changes fundamentally.

\medskip
\begin{quote}
Renormalization is not an operation that removes divergence.
\end{quote}
\medskip

Instead,

\medskip
\begin{quote}
Renormalization is the operation that converts \bcontact{}
into finite internal representation.
\end{quote}
\medskip

\subsection*{Two Layers of Physical Quantities}

In quantum field theory,
quantities appearing in calculations often diverge.

However,
observable quantities remain finite.

Renormalization connects these two layers.

Quantities containing divergences
are absorbed into the definition of observables,
and predictions emerge as finite values.

Conventionally,
this is regarded as a technical procedure.

But from the perspective of the \DH{},
its structure becomes clear.

\medskip
\begin{quote}
Theory cannot directly handle the information of \bcontact{}.
\end{quote}
\medskip

Divergence is the direct expression of that contact.
But in that form,
it cannot be consistently treated within the theory.

Thus, it is transformed.

\medskip
\begin{quote}
The divergence itself does not disappear.\\
Only its influence is retained in a transformed form.
\end{quote}
\medskip

\subsection*{Renormalization as Mapping}

\medskip
\begin{quote}
Renormalization is not elimination.\\
It is mapping.
\end{quote}
\medskip

More precisely,
it is a projection of \bcontact{} into a form
that can be handled within the theory.

\medskip
\begin{quote}
It is not removal,
but projection into a lower-resolution representation.
\end{quote}
\medskip

This explains why renormalization succeeds.

\medskip
\begin{quote}
Renormalization succeeds because it preserves the effects of \bcontact{}
while maintaining internal consistency of the theory.
\end{quote}
\medskip

\subsection*{Scale Dependence and Effective Theory}

Several known features of renormalization
follow naturally from this interpretation.

Scale dependence appears
because \bcontact{} depends on the scale
at which the theory is defined.

\medskip
\begin{quote}
Different scales correspond to different modes of \bcontact{}.
\end{quote}
\medskip

Effective field theory works
because the influence of high-energy structure
is encoded in low-energy observables.

\medskip
\begin{quote}
A physical quantity is a quantity that encodes
the history of \bcontact{}.
\end{quote}
\medskip

\subsection*{Structural Necessity}

Thus,
renormalization is not merely a correction procedure.

It is the fundamental mechanism by which a theory
remains valid within its own limits.

\medskip
\begin{quote}
Divergence indicates contact;\\
renormalization makes that contact visible.
\end{quote}

\medskip
\begin{quote}
Divergence and renormalization are not opposing operations,\\
but two aspects of a single structural process.
\end{quote}
\medskip

The next section, based on this perspective,
shows how major physical theories are placed
in relation to the \DH{}---
not by their objects, but by their mode of approach.
